\clearpage
\begin{landscape}
\begingroup
\fontsize{10}{13}\selectfont
\fancyhead[L]{\small NACHLASS. STÖRUNGEN DER PALLAS}
\renewcommand{\headrulewidth}{0pt}
\pdfbookmark[2]{Allgemeine Störungen, Erste Rechnung, [8]}{pallas-V-fragment8}
\begin{center}[8.]\end{center}
\par\noindent\hspace*{\parindent}[Zur Ermittelung der Störungen des Perihels wurde die Grösse \(\frac{\mathrm{d}\tilde\omega}{n\,\mathrm{d}t}-(1-\cos i)\frac{\mathrm{d}\Omega}{n\,\mathrm{d}t}\) nach der Gleichung 5, Artikel [1], berechnet, nach deren Integration die Störungen des Perihels ohne Weiteres gefunden werden können, sobald diejenigen der Länge des Knotens bekannt sind.
\par\noindent\hspace*{\parindent}Mit Hülfe der in der Tafel auf Seite 512--513 gegebenen Werthe von \(T\) und \(V\) fanden sich die Ausdrücke \(\frac{p}{e}\cos(v-\tilde\omega)\mathord{.}\,T\) und \(\frac{1}{e}\{2+e\cos(v-\tilde\omega)\}\sin(v-\tilde\omega)\mathord{.}\,V\), deren Summe alsdann nach der in den vorstehenden Artikeln angewandten Methode nach den Vielfachen von \(M\) entwickelt werde; das Resultat der Ausmultiplicirung soll der Raumersparniss wegen nicht angegeben werden.
\par\noindent\hspace*{\parindent}Die Integration dieses Ausdrucks geschah in der folgenden Weise: Sei ein Glied desselben \(b\cos z+b'\sin z\), wo \(z\), wie im vorigen Artikel, von der Form \(m'M'+mM\) ist, so lässt sich dieses Glied zusammenziehen in die Form
\[
q\sin(z+Q),
\]
\noindent wo
\[
\begin{aligned}
q\sin Q&=b\\
q\cos Q&=b'.
\end{aligned}
\]
\par\noindent\hspace*{\parindent}Das entsprechende Störungsglied wird dann
\[
-\frac{q}{m'\frac{n'}{n}+m}\cos(z+Q).
\]
\par\noindent\hspace*{\parindent}Für das constante Glied (\(m'=m=0\)) verschwinden \(z\) und \(b'\), \(Q\) wird gleich \(90^0\) und die jährliche Änderung von \(\tilde\omega-(1-\cos i)\Omega\) wird hienach gleich
\[
bn\mathord{.}\,365{,}25.
\]
\par\noindent\hspace*{\parindent}Als Beispiel diene wieder die Berechnung der grossen Gleichung, welche vom Argument \(5M'-2M\) abhängt:
\begingroup
\newcommand{\VEightPoint}{\mathord{.}\mkern5mu}
\newcommand{\VEightFive}{\VEightPoint\VEightPoint\VEightPoint\VEightPoint\VEightPoint}
\newcommand{\VEightComma}{\mathord{,}\mkern5mu}
\begin{center}
\begin{tabular}{@{}l@{\quad}l@{}}
\(\log b\;\VEightFive\VEightFive\VEightFive\VEightPoint\VEightPoint1{,}13532\)&\\
\(\log b'\;\VEightFive\VEightFive\VEightPoint\VEightPoint\VEightPoint\VEightPoint\VEightComma\VEightPoint\VEightPoint1{,}50678\)&\\
\multicolumn{1}{r}{\(Q=23^0\,1'\,58''\)}&\\
\(\log\cos Q\;\VEightFive\VEightFive\VEightPoint\VEightPoint\VEightPoint9{,}96392\)&also das Störungsglied:\\
\(\log q\;\VEightFive\VEightFive\VEightFive\VEightPoint\VEightPoint1{,}54286\)&\(+621''{,}16\cos(5M'-2M+23^0\,1'\,58'').\)\\
\(\log\left(5\frac{n'}{n}-2\right)\;\VEightFive\VEightPoint\VEightPoint\VEightPoint8{,}74966_n\)&\\
\cline{1-1}
\multicolumn{1}{r}{\(2{,}79320_n\)}&
\end{tabular}
\end{center}
\endgroup
\par\noindent\hspace*{\parindent}Die Störungen des Ausdrucks \(\tilde\omega-(1-\cos i)\Omega\) ergaben sich hienach, wie folgt:
\begingroup
\renewcommand{\arraystretch}{1.15}
\begin{longtable}{@{}r@{}l@{}}
\SourceLeaf{526}\(\delta\{\tilde\omega-(1-\cos i)\Omega\}={}\)&\(-54''{,}77\cos(M+174^0\,27'\,25'')\)\\
&\(-11{,}64\cos(2M+212\;29\;21)\)\\
&\(-5{,}05\cos(3M+352\;39\;52)\)\\
&\(-1{,}69\cos(4M+36\;1\;51)\)\\[3pt]
&\(+1{,}34\cos(M'-5M+48^0\,26'\,30'')\)\\
&\(+2{,}45\cos(M'-4M+84\;51\;45)\)\\
&\(+10{,}54\cos(M'-3M+244\;35\;58)\)\\
&\(+26{,}58\cos(M'-2M+282\;40\;0)\)\\
&\(+55{,}19\cos(M'-M+104\;14\;55)\)\\
&\(-66{,}05\cos(M'+36\;53\;29)\)\\
&\(-14{,}34\cos(M'+M+180\;13\;59)\)\\
&\(-3{,}48\cos(M'+2M+312\;11\;14)\)\\
&\(-2{,}02\cos(M'+3M+1\;43\;35)\)\\[3pt]
&\(+1{,}37\cos(2M'-6M+272^0\,49'\,40'')\)\\
&\(+1{,}72\cos(2M'-5M+240\;18\;29)\)\\
&\(+11{,}37\cos(2M'-4M+189\;15\;30)\)\\
&\(+26{,}56\cos(2M'-3M+218\;36\;16)\)\\
&\(+84{,}09\cos(2M'-2M+32\;14\;36)\)\\
&\(+785{,}02\cos(2M'-M+48\;20\;55)\)\\
&\(-84{,}04\cos(2M'+217\;34\;48)\)\\
&\(-12{,}77\cos(2M'+M+273\;39\;33)\)\\
&\(-1{,}80\cos(2M'+2M+323\;34\;0)\)\\[3pt]
&\(+1{,}03\cos(3M'-6M+99^0\,3'\,52'')\)\\
&\(+5{,}78\cos(3M'-5M+79\;49\;57)\)\\
&\(+6{,}69\cos(3M'-4M+137\;58\;58)\)\\
&\(+47{,}68\cos(3M'-3M+281\;13\;38)\)\\
&\(+92{,}48\cos(3M'-2M+321\;7\;28)\)\\
&\(-541{,}66\cos(3M'-M+98\;47\;17)\)\\
&\(-32{,}02\cos(3M'+213\;35\;51)\)\\
&\(-4{,}69\cos(3M'+M+314\;16\;0'')\)\\[3pt]
&\(+2{,}49\cos(4M'-6M+330^0\,22'\,10'')\)\\
&\(+2{,}77\cos(4M'-5M+96\;34\;55)\)\\
&\(+22{,}32\cos(4M'-4M+177\;8\;26)\)\\
&\(+25{,}60\cos(4M'-3M+253\;34\;38)\)\\
&\(+147{,}03\cos(4M'-2M+1\;37\;26)\)\\
&\(-77{,}67\cos(4M'-M+122\;10\;56)\)\\
&\(-12{,}11\cos(4M'+238\;41\;0)\)\\
&\(-2{,}25\cos(4M'+M+337\;10\;24)\)\\[3pt]
&\(+2''{,}24\cos(5M'-6M+19^0\,59'\,36'')\)\\
&\(+9{,}50\cos(5M'-5M+77\;6\;26)\)\\
&\(+12{,}99\cos(5M'-4M+184\;17\;47)\)\\
&\(+37{,}79\cos(5M'-3M+269\;36\;8)\)\\
&\(+621{,}16\cos(5M'-2M+23\;1\;58)\)\\
&\(-23{,}34\cos(5M'-M+148\;15\;51)\)\\
&\(-4{,}74\cos(5M'+260\;46\;22)\)\\
&\(-1{,}13\cos(5M'+M+358\;20\;48)\)\\[3pt]
&\(+1{,}45\cos(6M'-7M+283^0\,15'\,52'')\)\\
&\(+3{,}82\cos(6M'-6M+342\;59\;55)\)\\
&\(+7{,}14\cos(6M'-5M+98\;35\;24)\)\\
&\(+13{,}90\cos(6M'-4M+183\;20\;50)\)\\
&\(+39{,}71\cos(6M'-3M+295\;55\;45)\)\\
&\(-60{,}82\cos(6M'-2M+55\;16\;42)\)\\
&\(-8{,}24\cos(6M'-M+172\;47\;59)\)\\
&\(-2{,}00\cos(6M'+283\;39\;49)\)\\[3pt]
&\(+1{,}56\cos(7M'-7M+256^0\,11'\;4'')\)\\
&\(+3{,}76\cos(7M'-6M+6\;32\;34)\)\\
&\(+6{,}13\cos(7M'-5M+99\;54\;55)\)\\
&\(+13{,}58\cos(7M'-4M+204\;40\;35)\)\\
&\(+57{,}33\cos(7M'-3M+321\;44\;53)\)\\
&\(-14{,}86\cos(7M'-2M+80\;41\;55)\)\\
&\(-3{,}14\cos(7M'-M+194\;14\;3)\)\\[3pt]
&\(+2{,}05\cos(8M'-7M+276^0\,8'\,52'')\)\\
&\(+3{,}23\cos(8M'-6M+18\;25\;37)\)\\
&\(+5{,}94\cos(8M'-5M+119\;3\;17)\)\\
&\(+13{,}10\cos(8M'-4M+232\;33\;27)\)\\
&\(-84{,}68\cos(8M'-3M+351\;57\;34)\)\\
&\(-5{,}21\cos(8M'-2M+113\;55\;49)\)\\
&\(-1{,}43\cos(8M'-M+235\;36\;29)\)\\[3pt]
&\(+1{,}57\cos(9M'-7M+287^0\,15'\,29'')\)\\
&\(+2{,}68\cos(9M'-6M+27\;34\;10)\)\\
&\(+5{,}07\cos(9M'-5M+137\;15\;53)\)\\
&\(+13{,}81\cos(9M'-4M+255\;7\;55)\)\\
&\(-9{,}80\cos(9M'-3M+12\;27\;32)\)\\
&\(-1{,}90\cos(9M'-2M+129\;41\;51)\)\\[3pt]
&\(+1{,}23\cos(10M'-7M+287^0\,35'\;5'')\)\\
&\(+2{,}34\cos(10M'-6M+46\;5\;0)\)\\
\SourceLeaf{527}&\(+4''{,}46\cos(10M'-5M+160^0\,47'\;3'')\)\\
&\(+35{,}47\cos(10M'-4M+278\;40\;22)\)\\
&\(-2{,}94\cos(10M'-3M+36\;23\;44)\)\\[3pt]
&\(+1{,}13\cos(11M'-7M+316^0\,32'\,48'')\)\\
&\(+1{,}95\cos(11M'-6M+69\;1\;13)\)\\
&\(+4{,}17\cos(11M'-5M+186\;28\;5)\)\\
&\(-8{,}01\cos(11M'-4M+304^0\,25'\,18'')\)\\
&\(-1{,}06\cos(11M'-3M+63\;10\;13)\)
\end{longtable}
\begin{center}
\begin{tabular}{rl}
Säcular: jährlich&\(-4''{,}787\)\\
täglich&\(-0{,}01310\)\\[4pt]
\multicolumn{2}{l}{Säcularänderung von \(\tilde\omega\):}\\[4pt]
jährlich&\(-11''{,}015\)\\
täglich&\(-0''{,}03015\).]
\end{tabular}
\end{center}
\endgroup
\clearpage
\endgroup
\end{landscape}
